\section{Discusión}

\subsection{Analogía limitada con sistemas de múltiples escalas}

La dinámica entre capas puede visualizarse, de forma heurística y limitada, mediante la analogía de un huracán maduro. Las intensificaciones locales (celdas convectivas intensas) corresponden a presentes locales intensificados. La organización entre escalas (bandas, eyewall, circulación general) corresponde a la reducción de fragmentación y al aumento de coherencia entre presentes locales. La inercia del huracán maduro ilustra el efecto de filtrado probabilístico de las capas superiores frente a intensificaciones locales fragmentadas. La desorganización (cizalladura, interacción con tierra) corresponde al aumento de antisincronización.

La analogía es útil para transmitir que la intensificación local no equivale automáticamente a cambio estructural, y que la organización entre escalas es condición necesaria para que las intensificaciones locales tengan consecuencias globales. Sin embargo, presenta limitaciones importantes: en el huracán la organización tiene un fuerte componente físico y energético, mientras que en el presente la coherencia es fundamentalmente relacional. Además, el huracán posee un ciclo de vida claro, a diferencia del presente de capas superiores, cuyo posible patrón de consolidación, madurez o eventual fragmentación permanece por determinar.

\subsection{Preguntas abiertas}

\begin{enumerate}
    \item ¿En qué condiciones precisas una intensificación del presente local logra superar el efecto de filtrado probabilístico ascendente de las capas superiores y generar reorganización estructural efectiva?
    
    \item ¿Existe un análogo de ``ciclo de vida'' del presente de capas superiores (consolidación, madurez, posible decadencia o aumento de fragmentación)?
    
    \item ¿Cómo se comporta la dinámica de fragmentación y filtrado probabilístico en sistemas con acoplamiento fuertemente asimétrico (un módulo con peso dominante)?
    
    \item ¿Es posible detectar empíricamente casos de retrocausalidad local, entendida como reversión hacia configuraciones estructurales anteriores mediada por alta antisincronización sostenida?
    
    \item ¿Qué rol preciso juega la memoria activa (trapping + hiper-persistencia) en el efecto de filtrado probabilístico descendente frente a intensificaciones locales fragmentadas?
    
    \item ¿Cómo se articula esta visión estratificada del presente con concepciones filosóficas del tiempo de la esencia, en particular con el proyecto antropológico de Leonardo Polo (véase \citep{polo_antropologia_trascendental,padilla_polo_dialogue_2026})?
\end{enumerate}